\subsubsection{真实输入核的一参数加权 $\zeta(z,u)$（输出位势）}\label{app:real-input-40-zeta-u}
只对输出位势 $\varphi_e=e$ 加权，令 $u=e^{\theta_e}$，记相应的加权矩阵为 $M(u)$，并定义
\[
\zeta(z,u)=\frac{1}{\Delta(z,u)},\qquad \Delta(z,u):=\det(I-zM(u)).
\]
在真实输入 40 状态核上，$\Delta(z,u)$ 仍可压缩为低阶多项式并显式因子化：
\[
\Delta(z,u)=(1-uz^2)\Bigl(
1-z-6uz^2+2uz^3+(9u^2-u)z^4+(u-3u^2)z^5-(u^3+2u^2)z^6
+(4u^3-3u^2)z^7+(u^3-u^4)z^8
\Bigr).
\]
其中因子 $1-uz^2$ 并非“偶然代数分解”：它对应一条可在 primitive 层面被\textbf{精确定位}的原子闭轨道，并由此给出 Abel 有限部常数的精确剥离。

\begin{proposition}[原子 \texorpdfstring{$(2,1)$}{(2,1)}-prime orbit 与有限部常数的精确剥离]\label{prop:atomic-2-orbit-split-logM}
对加权矩阵族 $M(u)$，令每条 primitive 周期轨道 $\tau$ 的长度为 $|\tau|$，并记其一圈累计输出 $1$ 的次数为
\[
E(\tau):=\sum_{e\in\tau} e\in\ZZ_{\ge 0}.
\]
则有且仅有一条 primitive 轨道 $\tau_2$ 满足
\[
|\tau_2|=2,\qquad E(\tau_2)=1,
\]
并且该轨道在加权 zeta 的 Euler 乘积中贡献出因子 $(1-uz^2)^{-1}$。

进一步，令未加权 primitive 生成函数
\(
P(z)=\sum_{n\ge 1}p_n z^n
\)（$p_n$ 为未加权 primitive 轨道数，按轨道计），并定义“去掉原子 $2$-轨道后的核心部分”
\[
P_{\mathrm{core}}(z):=P(z)-z^2.
\]
则在 $z\to z_\star=\lambda^{-1}$ 时
\[
P_{\mathrm{core}}(z)=-\log(1-\lambda z)+\log\mathfrak{M}_{\mathrm{core}}+o(1),
\]
并满足常数的精确分解
\[
\boxed{\ \log\mathfrak{M}=\log\mathfrak{M}_{\mathrm{core}}+z_\star^2\ }.
\]
在本核上 $z_\star=\varphi^{-2}$，故
\[
z_\star^2=\varphi^{-4}=\frac{7-3\sqrt5}{2}\approx 0.145898033750315.
\]
因此由上文数值 $\log\mathfrak{M}\approx -0.254312661901716$ 可得
\[
\log\mathfrak{M}_{\mathrm{core}}\approx -0.400210695652031,
\qquad
\mathsf{M}_{\mathrm{core}}:=\gamma+\log\mathfrak{M}_{\mathrm{core}}\approx 0.177004969249501.
\]
\end{proposition}

\begin{proof}
对带权 SFT，闭轨道 Euler 乘积给出
\[
\zeta(z,u)=\prod_{\tau\ \mathrm{primitive}}\left(1-u^{E(\tau)}z^{|\tau|}\right)^{-1}
\qquad\text{（参见 \cite{BrinStuck2002,ParryPollicott1990Zeta} 的同型推导）}.
\]
与 $\Delta(z,u)=\det(I-zM(u))$ 的显式分解相比，因子 $(1-uz^2)^{-1}$ 只能由满足 $(|\tau|,E(\tau))=(2,1)$ 的 primitive 轨道贡献。
若存在两条不同的此类轨道，则该因子在乘积中应出现二次幂；与行列式分解矛盾，故此类轨道存在且唯一。

另一方面，$u=1$ 时该轨道在未加权 primitive 生成函数中贡献 $z^{|\tau_2|}=z^2$，从而
$P(z)=z^2+P_{\mathrm{core}}(z)$。由于 $z^2$ 在 $z=z_\star$ 处解析，故 Abel 有限部常数相差恰为
$z_\star^2$，即 $\log\mathfrak{M}=\log\mathfrak{M}_{\mathrm{core}}+z_\star^2$。
\end{proof}

\begin{proposition}[输出压力的冻结端点与 \texorpdfstring{$\sqrt u$}{sqrt(u)} 标度]\label{prop:real-input-40-pressure-freezing}
令 $u=e^{\theta}$，记
\[
\lambda(u):=\rho(M(u)),\qquad P(\theta):=\log\lambda(e^{\theta}).
\]
把 $\Delta(z,u)=(1-uz^2)F(z,u)$ 中的 $F$ 写为 $\lambda$-多项式
\[
G(\lambda,u):=\lambda^8\,F(\lambda^{-1},u)
=\lambda^8-\lambda^7-6u\lambda^6+2u\lambda^5+(9u^2-u)\lambda^4+(u-3u^2)\lambda^3-(u^3+2u^2)\lambda^2
+(4u^3-3u^2)\lambda+(u^3-u^4).
\]
则 $\lambda(u)$ 为 $G(\lambda,u)=0$ 的最大正实根，并且：
\begin{enumerate}
  \item \textbf{$u\to 0$（$\theta\to-\infty$）冻结：}有 $\lambda(u)\to 1$，并且在 $u=0$ 的正规微扰给出
  \[
  \lambda(u)=1+4u-5u^2-115u^3+2317u^4+O(u^5),
  \]
  从而
  \(
  P(\theta)=4e^{\theta}+O(e^{2\theta})
  \)。
  \item \textbf{$u\to\infty$（$\theta\to+\infty$）饱和：}存在常数 $c>0,d\in\RR$ 使得
  \[
  \lambda(u)=c\sqrt u+d+O(u^{-1/2}),
  \]
  其中 $y:=c^2$ 为四次方程
  \[
  y^4-6y^3+9y^2-y-1=0
  \]
  的最大正根；并且
  \[
  d=\frac{y^3-2y^2+3y-4}{8y^3-36y^2+36y-2}\Bigg|_{y=c^2}.
  \]
  数值上 $c\approx 1.89635299351377,\ d\approx 0.807943317323508$，从而
  \[
  P(\theta)=\frac{\theta}{2}+\log c+\frac{d}{c}e^{-\theta/2}+O(e^{-\theta}).
  \]
\end{enumerate}
因此
\[
\lim_{\theta\to-\infty}P'(\theta)=0,\qquad \lim_{\theta\to+\infty}P'(\theta)=\frac12,
\]
即输出 $1$ 的典型密度在极端偏置下冻结在区间 $[0,\tfrac12]$。
\end{proposition}

\begin{corollary}[端点大偏差代价]\label{cor:real-input-40-rate-endpoints}
设 $I(\alpha)$ 为输出位势 $\varphi_e=e$ 的（归一化）大偏差速率函数（以 $P(\theta)$ 的 Legendre 形式定义）。则端点满足
\[
I(0)=P(0)=\log\lambda(1)=\log(\varphi^2)\approx 0.962423650119207,
\]
\[
I(1/2)=P(0)-\log c=\log\!\left(\frac{\varphi^2}{c}\right)\approx 0.322491085601762.
\]
\end{corollary}
\begin{proposition}[极端偏置的有限部常数极限]\label{prop:real-input-40-logM-infty}
令 $z_\star(u)=\lambda(u)^{-1}$，并定义
\[
b:=\lim_{u\to\infty}u\,z_\star(u)^2=\frac{1}{c^2}.
\]
则 $b$ 为方程
\[
1-6b+9b^2-b^3-b^4=0
\]
的最大正根，数值为 $b\approx 0.27807480214113$。记
\[
C(u):=\lim_{z\to z_\star(u)}(1-\lambda(u)z)\,\zeta(z,u),
\]
则存在极限
\[
C_\infty:=\lim_{u\to\infty}C(u)=\frac{1}{2b(1-b)(6-18b+3b^2+4b^3)}\approx 1.89745149363.
\]
此外 Abel 有限部常数满足极限
\[
\log\mathfrak{M}_\infty:=\lim_{u\to\infty}\log\mathfrak{M}(u),
\]
并具有快速收敛的 Möbius--Euler 表达式
\[
\log\mathfrak{M}_\infty
=\log C_\infty
-\sum_{m\ge 2}\frac{\mu(m)}{m}\log\Bigl[(1-b^m)\bigl(1-6b^m+9b^{2m}-b^{3m}-b^{4m}\bigr)\Bigr].
\]
数值上
\[
\log\mathfrak{M}_\infty\approx 0.282290255262,\qquad
\mathfrak{M}_\infty\approx 1.32616359009.
\]
\end{proposition}
该闭式对象的主极点给出输出位势下的一维压力曲线 $P(\theta_e)=\log\rho(M(e^{\theta_e}))$，并可直接用于谱分支切换/相变检查（因式 $1-uz^2$ 对应一条可追踪的谱分支）。
数值核验与复现见脚本 \texttt{scripts/exp\_sync\_kernel\_real\_input\_40\_zeta\_u.py}。
